% =========================================================
% The Ordered Field Structure of Q
% =========================================================

\section{The Ordered Field Structure of $\mathbb{Q}$}
\label{sec:ordered-field-structure-q}

Once the order has been defined, the main task is to prove that it is compatible with the rational field operations.

\begin{theorem}[$\mathbb{Q}$ Is an Ordered Field]
\label{thm:q-is-ordered-field}
With its usual order, the system
\[
(\mathbb{Q},+,\cdot,0,1,<)
\]
is an ordered field.
\end{theorem}

\begin{remark*}[Proof]
\hyperref[prf:q-is-ordered-field]{Go to proof.}
\end{remark*}

\begin{remark}
This theorem records the compatibility of the rational order with the field
operations. It is this compatibility that makes inequalities in $\mathbb{Q}$
behave as expected under addition, subtraction, multiplication, and division
by positive elements.
\end{remark}

\begin{theorem}[Squares Are Nonnegative]
\label{thm:squares-nonnegative}
In any ordered field $F$, for every $a\in F$,
\[
a^2\ge 0.
\]
\end{theorem}

\begin{remark*}[Proof]
\hyperref[prf:squares-nonnegative]{Go to proof.}
\end{remark*}

\begin{theorem}[Positive Elements Have Positive Inverses]
\label{thm:inverse-positive}
In any ordered field, if $a>0$, then $a^{-1}>0$.
\end{theorem}

\begin{remark*}[Proof]
\hyperref[prf:inverse-positive]{Go to proof.}
\end{remark*}

\begin{theorem}[Multiplication by a Positive Element Preserves Order]
\label{thm:multiply-positive}
In any ordered field, if $a<b$ and $c>0$, then
\[
ac<bc.
\]
\end{theorem}

\begin{remark*}[Proof]
\hyperref[prf:multiply-positive]{Go to proof.}
\end{remark*}

\begin{theorem}[Multiplication by a Negative Element Reverses Order]
\label{thm:multiply-negative}
In any ordered field, if $a<b$ and $c<0$, then
\[
ac>bc.
\]
\end{theorem}

\begin{remark*}[Proof]
\hyperref[prf:multiply-negative]{Go to proof.}
\end{remark*}

\begin{theorem}[$1$ Is Positive]
\label{thm:one-positive}
In any ordered field,
\[
1>0.
\]
\end{theorem}

\begin{remark*}[Proof]
\hyperref[prf:one-positive]{Go to proof.}
\end{remark*}

\begin{corollary}[Natural Numbers Are Positive]
\label{cor:n-positive}
In any ordered field, for every natural number $n$,
\[
n>0.
\]
\end{corollary}

\begin{remark*}[Proof]
\hyperref[prf:n-positive]{Go to proof.}
\end{remark*}

\begin{theorem}[Positivity of a Quotient via the Product $ab$]
\label{thm:positive-fraction-iff-positive-product}
For any $a,b\in F$ with $b\neq 0$,
\[
0<\frac{a}{b}
\quad\Longleftrightarrow\quad
0<a\cdot b.
\]
\end{theorem}

\begin{remark*}[Proof]
\hyperref[prf:positive-fraction-iff-positive-product]{Go to proof.}
\end{remark*}

\begin{theorem}[General Fraction Comparison in an Ordered Field]
\label{thm:fraction-order-comparison-general}
For any $a,b,c,d\in F$ with $b\neq 0$ and $d\neq 0$,
\[
\frac{a}{b}<\frac{c}{d}
\quad\Longleftrightarrow\quad
(a\cdot d)\cdot(b\cdot d)<(b\cdot c)\cdot(b\cdot d).
\]
\end{theorem}

\begin{remark*}[Proof]
\hyperref[prf:fraction-order-comparison-general]{Go to proof.}
\end{remark*}

\begin{theorem}[Fraction Comparison When the Denominator Product Is Positive]
\label{thm:fraction-order-comparison-positive-denominator-product}
For any $a,b,c,d\in F$, if $b\neq 0$, $d\neq 0$, and $0<b\cdot d$, then
\[
\frac{a}{b}<\frac{c}{d}
\quad\Longleftrightarrow\quad
a\cdot d<b\cdot c.
\]
\end{theorem}

\begin{remark*}[Proof]
\hyperref[prf:fraction-order-comparison-positive-denominator-product]{Go to proof.}
\end{remark*}

\begin{theorem}[Reciprocal Order for Positive Elements]
\label{thm:reciprocal-order-for-positive-elements}
If $0<d<b$, then
\[
0<\frac{1}{b}<\frac{1}{d}.
\]
Conversely, if
\[
0<\frac{1}{b}<\frac{1}{d},
\]
then $0<d<b$.
\end{theorem}

\begin{remark*}[Proof]
\hyperref[prf:reciprocal-order-for-positive-elements]{Go to proof.}
\end{remark*}

\begin{theorem}[Ordered Fields Are Densely Ordered]
\label{thm:ordered-fields-are-densely-ordered}
Every ordered field is densely ordered by its order relation.
\end{theorem}

\begin{remark*}[Proof]
\hyperref[prf:ordered-fields-are-densely-ordered]{Go to proof.}
\end{remark*}

\begin{remark}[Transition]
\label{rem:arch-v-complete}
The positivity of $1$ anchors the order structure and leads naturally to
the Archimedean property. That property explains why natural numbers
eventually outrun every fixed rational and why reciprocals like $1/n$
become arbitrarily small.
\end{remark}
